%Chargement des paquets
\usepackage{amsmath}
\usepackage{amsfonts}
\usepackage{amssymb}
\usepackage{enumerate}
\usepackage{mathtools}
\usepackage{bbm}
\usepackage{xparse, etoolbox}
\usepackage{enumerate}
\usepackage{enumitem}
\usepackage{mathabx}
\usepackage{babel}[french]

%environment
\usepackage{amsthm}
\usepackage{keytheorems}
\usepackage{hyperref} 

%Interface théorème
\newkeytheorem{lem}[
	name = Lemme
]

\newkeytheorem{prop}[
	name = Proposition
]

\newkeytheorem{thm}[
	name = Théorème
]

\newkeytheorem{coro}[
	name = Corollaire
]

\renewcommand*{\proofname}{Démonstration}

\theoremstyle{definition}
\newtheorem{defn}{Définition}

\theoremstyle{definition}
\newtheorem*{nota}{Notation}

\theoremstyle{definition}
\newtheorem*{limi}{Liminaire}

\theoremstyle{remark}
\newtheorem{rmq}{Remarque}

\theoremstyle{plain}
\newtheorem*{fait}{Fait}


%Ensembles de nombres
\newcommand{\N} {\mathbb{N}}
\newcommand{\Ne}{\N^\ast}
\newcommand{\Z} {\mathbb{Z}}
\newcommand{\Q} {\mathbb{Q}}
\newcommand{\R} {\mathbb{R}}
\newcommand{\Rb}{\overline{\mathbb{R}}}
\newcommand{\Rp}{\R_+}
\newcommand{\Rm}{\R_-}
\newcommand{\K} {\mathbb{K}}
\newcommand{\Ket} {\mathbb{K^{\ast}}}
\newcommand{\Cx}{\mathbb{C}}

%Ensembles, tribus
\newcommand{\A}{\mathbb{A}}
\renewcommand{\P}{\mathcal{P}}
\newcommand{\T}{\mathcal{T}}
\newcommand{\F}{\mathcal{F}}
\newcommand{\C} {\mathcal{C}}
\newcommand{\ouv}{\mathcal{O}}
\newcommand{\B}{\mathcal{B}}
\newcommand{\M}{\mathcal{M}}
\newcommand{\etag}{\mathcal{E}}
\newcommand{\etagOp}{\etag^{+}(\Omega)}
\newcommand{\etagO}{\etag(\Omega)}
%%Lp avant quotientage
\newcommand{\Lic}{\mathcal{L}^1}
\newcommand{\Lpc}{\mathcal{L}^p}
\newcommand{\Linfc}{\mathcal{L}^{\infty}}
\newcommand{\Lifull}{\Lic(\Omega, \B, m)}
\newcommand{\Lpfull}{\Lpc(\Omega, \B, m)}
\newcommand{\Linffull}{\Linfc(\Omega, \B, m)}
%%Lp après quotientage
\newcommand{\Li}{L^1}
\newcommand{\Ld}{L^2}
\newcommand{\Lp}{L^p}
\newcommand{\Linf}{L^{\infty}}
\newcommand{\Listd}{\Li(\Omega, m)} 
\newcommand{\Ldstd}{\Ld(\Omega, m)} 
\newcommand{\Lpstd}{\Lp(\Omega, m)} 

%Quelques opérateurs
\DeclareMathOperator{\codim}{codim}
\DeclareMathOperator{\rg}{rg}
\DeclareMathOperator{\tr}{tr}
\DeclareMathOperator{\vect}{Vect}
\DeclareMathOperator{\ima}{Im}
\DeclareMathOperator{\id}{Id}
\DeclareMathOperator{\sign}{sign}
\DeclareMathOperator{\ext}{ext}
\newcommand{\ind}{\mathbbm{1}}
\newcommand*{\spr}[2]{\left(\,#1\,\middle|\,#2\,\right)}
\newcommand{\interior}[1]{%
  {\kern0pt#1}^{\mathrm{o}}%
}
\DeclareMathOperator{\card}{card}
\DeclareMathOperator{\ppcm}{ppcm}

%Commandes de pure flemme
\newcommand{\veps}{\varepsilon}
\newcommand{\intab}{\int_{a}^{b}}
\newcommand{\intR}{\int_{-\infty}^{\infty}}
\newcommand{\intinf}{\int_{- \infty}^{+ \infty}}
\newcommand{\equi}{\Leftrightarrow}
\newcommand{\Kev}{$\K$-espace vectoriel }
\newcommand{\Kevs}{$\K$-espaces vectoriels }
\newcommand{\Rev}{$\R$-espace vectoriel }
\newcommand{\Revs}{$\R$-espaces vectoriels }
\newcommand{\Cev}{$\Cx$-espace vectoriel }
\newcommand{\Cevs}{$\Cx$-espaces vectoriels }
\newcommand{\evn}{(E, \norm{.})}
\newcommand{\interf}[2]{[#1,#2]}
\newcommand{\limb}{\lim\limits_}
\newcommand{\lebn}{\lambda_n}
\newcommand{\pp}{\text{~presque partout}}
\newcommand{\derivp}[2]{\frac{\partial #1}{\partial #2}}
\newcommand{\famiu}{(u_i)_{i \in I}}
\newcommand{\sev}{\text{~s.e.v.}}
\newcommand{\Mnp}{M_{n,p}(\K)}
\newcommand{\Mn}{M_{n}(\K)}
\newcommand{\tq}{\text{~t.q.~}}
\newcommand{\mq}{~montrer~que~}
\newcommand{\et}{\quad \text{et} \quad}
\newcommand{\ou}{\text{~ou~}}
\newcommand{\Kx}{\K[X]}


%Commandes de gars chelou
\renewcommand{\subset}{\subseteq}

%Commande restriction, ty duckduckgo
\def\restr#1#2{\mathchoice
              {\setbox1\hbox{${\displaystyle #1}_{\scriptstyle #2}$}
              \restraux{#1}{#2}}
              {\setbox1\hbox{${\textstyle #1}_{\scriptstyle #2}$}
              \restraux{#1}{#2}}
              {\setbox1\hbox{${\scriptstyle #1}_{\scriptscriptstyle #2}$}
              \restraux{#1}{#2}}
              {\setbox1\hbox{${\scriptscriptstyle #1}_{\scriptscriptstyle #2}$}
              \restraux{#1}{#2}}}
\def\restraux#1#2{{#1\,\smash{\vrule height .8\ht1 depth .85\dp1}}_{\,#2}} 

%Quelques normes
\DeclarePairedDelimiter\abs{\lvert}{\rvert}
\DeclarePairedDelimiter\labs{\bigl\lvert}{\bigr\rvert}
\DeclarePairedDelimiter\norm{\lVert}{\rVert}
\DeclarePairedDelimiterXPP\normi[1]{}\lVert\rVert{_1}{\ifblank{#1}{\:\cdot\:}{#1}}
\DeclarePairedDelimiterXPP\normd[1]{}\lVert\rVert{_2}{\ifblank{#1}{\:\cdot\:}{#1}}
\DeclarePairedDelimiterXPP\normp[1]{}\lVert\rVert{_p}{\ifblank{#1}{\:\cdot\:}{#1}}
\DeclarePairedDelimiterXPP\normq[1]{}\lVert\rVert{_q}{\ifblank{#1}{\:\cdot\:}{#1}}
\DeclarePairedDelimiterXPP\norminf[1]{}\lVert\rVert{_\infty}{\ifblank{#1}{\:\cdot\:}{#1}}

%Bracket en angle
\DeclarePairedDelimiter\angl{\langle}{\rangle}

%Set, parenthèses
\newcommand{\set}[1]{\{ #1 \}}
\newcommand{\bigset}[1]{\bigl\{ #1 \bigr\}}
\newcommand{\Bigset}[1]{\Bigl\{ #1 \Bigr\}}
\newcommand{\bigpar}[1]{\bigl( #1 \bigr)}
\newcommand{\Bigpar}[1]{\Bigl( #1 \Bigr)}

%Opitmisation des opérateurs de comparaison
\DeclarePairedDelimiter\tio{o (}{)}
\DeclarePairedDelimiter\gro{O (}{)}

\newcommand{\Ee}{E^{\ast}}

\pagestyle{fancy}

\fancyhead[C]{Théorème de projection sur un convexe}
\fancyhead[R]{}

\begin{document}

\underline{Sources :} 
\begin{itemize}
\item Hassan - Topologie - page 355 (353 pour caractérisation) ;
\item Hirsch Lacombe - analyse - page 91 ;
\end{itemize}

\underline{Leçons :}
\begin{itemize}
\item  205 : Espaces complets. Exemples et applications.
\item  213 : Espaces de Hilbert. Exemples d’applications.
\item  208 : Espaces vectoriels normés, applications linéaires continues. Exemples.
\item  219 : Extremums : existence, caractérisation, recherche. Exemples et applications.
\item  253 : Utilisation de la notion de convexité en analyse.
\end{itemize}

\begin{nota}
Le corps $\K$ est $\R$ ou $\Cx$ et $(E, \angl{,})$ est un $\K$-espace préhilbertien.
\end{nota}

\begin{lem}[Identité du parallélogramme]
Pour tous $x, y \in E$, on a :
\[ \norm{x+y}^2 + \norm{x-y}^2 = 2( \norm{x}^2 + \norm{y}^2) . \]
\end{lem}

\begin{proof}
Par calcul direct :
\begin{align*}
\angl{x+y,x+y} + \angl{x-y,x-y} & = \angl{x,x} + \angl{x,y} + \angl{y,x} + \angl{y,y} + \angl{x,x} - \angl{x,y} - \angl{y,x} + \angl{y,y} \\
& = 2\angl{x,x} + 2\angl{y,y} 
\end{align*}
\end{proof}

\begin{thm}
Soit $A$ une partie non vide convexe de $E$. 
On suppose de plus que $A$ est complète lorsque munie de la distance induite par la norme sur $E$ (ce qui est le cas pour $E$ Hilbertien
et $A$ une partie non vide convexe et fermée de $E$).
Pour tout $x \in E$, il existe un unique point $a \in A$ tel que :
\[ \norm{x-a} = d(x,A) . \]
Ce point est appelé le projeté de $x$ sur $A$ et est noté $P_A(x)$.
Ce projeté est caractérisée par :
\begin{equation}
\forall z \in A, \quad \Re{\angl{x-a,z-a}} \leq 0 . \label{eq:1}
\end{equation}
\end{thm}

\begin{proof}
Il est clair que si $E$ est complet et que si $A$ est fermé, alors, toute suite de Cauchy d'éléments de $A$ converge dans $A$. \\

Considérons $x \in E$ et prenons $(a_n)_\N$ une suite de $A$ telle que $\norm{x-a_n} \to d(x,A)$ lorsque $n \to \infty$. 
On remarque qu'une telle suite existe car, par définition de la borne inférieure :
\[ \forall n \in \N, \exists a_n \in A \tq \norm{a_n-x} \leq d(x,A) + 1/2^n . \]
Considérons $n, m \in \N$, on applique l'identité du parallélogramme aux points $x-a_n$ et $x-a_m$ : 
\[ \norm{2x- a_n - a_m}^2 + \norm{a_n - a_m}^2 = 2(\norm{x-a_n}^2 + \norm{x-a_m}^2) , \]
d'où on déduit :
\[ \norm{a_n-a_m}^2 = 2(\norm{x-a_n}^2 + \norm{x-a_m}^2) - 4 \norm{x - \dfrac{a-n+a_m}{2}}^2 . \]
Par convexité de $A$, on a $\dfrac{a-n+a_m}{2} \in A$ et donc $\norm{x - \dfrac{a-n+a_m}{2} }^2 \geq d(x,A)$, soit :
\begin{equation}
\norm{a_n-a_m}^2 \leq 2(\norm{x-a_n}^2 + \norm{x-a_m}^2) - 4 d(x,A) , \label{eq:2} 
\end{equation}
et clairement, comme $\norm{x-a_n} \to 0$, on déduit que la suite $(a_n)_\N$ est de Cauchy dans $A$, donc convergente vers $a \in A$.
Finalement, on a $d(x,A) = \norm{x-a}$. \\

Supposons qu'il existe deux tels points de $A$ notés $a$ et $b$, alors, selon l'équation (\ref{eq:2}), on a :
\[ \norm{a-b}^2 \leq 2(\norm{x-a}^2 + \norm{x-b}^2) - 4 d(x,A) = 0 , \]
soit $a=b$. \\

Il reste à montrer la caractérisation du projeté. Soit $z \in A$, par convexité le segment $[z,a]$ est inclus dans $A$ donc :
\[ \forall t \in [0,1], \quad \norm{x-[(1-t)a+tz]}^2 \geq \norm{x-a}^2 , \]
soit :
\begin{align*}
\norm{x-a+t(a-z)}^2 \geq \norm{x-a}^2 & \\
& \Rightarrow \norm{x-a}^2 + 2t \Re(\angl{x-a,a-z}) + t^2 \norm{a-z}^2 \geq \norm{x-a}^2 \\
& \Rightarrow \Re{\angl{x-a,z-a}} \leq \dfrac{t}2 \norm{a-z}^2 .
\end{align*}
Ainsi, en faisant tendre $t$ vers 0, on obtient l'inégalité recherchée. \\
Réciproquement, si $a \in A$ vérifie (\ref{eq:1}), alors, pour tout $z \in A$ : 
\begin{align*}
\norm{x-z}^2 & = \norm{x-a + a - z}^2 \\
& = \norm{x-a}^2 + \norm{a-z}^2 + 2 \Re(\angl{x-a, a-z}) \geq \norm{x-a}^2 ,
\end{align*}
autrement dit, $a = \inf{z \in C ; \norm{x-z} }$.
\end{proof}

\begin{thm}
Sous les conditions du théorème précédent l'application $P_C$ qui, à un point $x \in E$ associe son projeté $P_A(x)$ sur $A$ est 1-lipschitizenne.
\end{thm}

\begin{proof}
Soient $x, y \in E$ et notons $a = P_A(x), b=P_A(x')$ et $z=x-y-(a-b)$. On a :
\begin{align*}
\norm{x-y}^2 & = \norm{z + (a-b)}^2 \\
& = \norm{z}^2 + \norm{a-b}^2 + 2\Re{\angl{z, a-b}} . 
\end{align*}
D'autre part :
\begin{align*}
\Re{\angl{z, a-b}} & = \Re(\angl{x-a - (y-b), a-b}) \\
& = - \Re{\angl{x-a, b-a}} - \Re(\angl{y-b, a-b}) \geq 0 \qquad (\text{caractérisation du projeté}) .
\end{align*}
On a ainsi montré que $\norm{x-y} \geq \norm{a-b} = \norm{P_A(x) - P_A(y)}$, ce qui montre que l'application $P_A$ est 1-lipschitizienne,
en particulier continue. 
\end{proof}

Pour aller plus loin, on peut proposer le résultat suivant :

\begin{thm}
L'espace $E$ est la somme directe topologique de $A$ et $A^{\perp}$ et, pour tout $x \in E$, on a $a = P_A(x)+P_{A^{\perp}}(x)$.
\end{thm}

\end{document}